\chapter{Analyse du besoin}

\section{Description du processus de mesure souhaité}

Le client souhaite mettre en place une méthode de mesure accélérée de la réserve de marche d’un mouvement mécanique. 
Plutôt que de mesurer la durée complète de fonctionnement du mouvement, il est proposé de ne mesurer que les derniers tours du barillet, correspondant à la fin de la réserve de marche.

La séquence de mesure souhaitée est la suivante :

\begin{itemize}
    \item Le mouvement est placé par l’horloger dans un posage adaptable à différents mouvements (diamètre, position de la tige, présence d’aiguilles ou de cadran).
    \item La machine remonte le mouvement d’un nombre défini de tours afin de garantir un armage complet du barillet.
    \item La tige de remontoir est ensuite bloquée en rotation, laissant le mouvement fonctionner librement.
    \item La machine dévide artificiellement le barillet d’un nombre défini de tours, correspondant à une partie de la réserve de marche totale.
    \item La tige est à nouveau bloquée et le mouvement est laissé en fonctionnement jusqu’à son arrêt.
    \item Le moment de l’arrêt est détecté automatiquement afin de déduire la réserve de marche réduite.
\end{itemize}

\section{Contraintes et exigences}

La machine devra répondre aux contraintes suivantes :

\begin{itemize}
    \item Compatibilité avec différents calibres de mouvements mécaniques.
    \item Absence de dégradation du mouvement testé.
    \item Précision suffisante de la mesure de la réserve de marche réduite.
    \item Réduction significative du temps de mesure par rapport à une mesure classique.
    \item Utilisation simple et reproductible par un opérateur.
\end{itemize}

\section{Critères de performance}

Les critères suivants permettront d’évaluer la qualité de la solution développée :

\begin{itemize}
    \item Gain de temps obtenu par rapport à une mesure complète de la réserve de marche.
    \item Répétabilité des mesures.
    \item Précision sur la détection du moment d’arrêt du mouvement.
    \item Facilité d’adaptation à différents mouvements.
\end{itemize}